18个独立输入的90次主实验中，G、LS、SLS、MIP和LPR分别获得14、14、12、17、12例可行方案。全部已接受方案通过独立的坐标、集合和调用核对。表\ref{tab:success}保留各方法失败次数，不因目标值缺失而缩小成功率分母。

LS与G在14个共同可行输入上比较，前者取得14优、0平、0劣，平均遗漏差为-9.5。

SLS与LS在8个共同可行输入上比较，前者取得3优、1平、4劣，平均遗漏差为-2.1。

LPR与LS在11个共同可行输入上比较，前者取得1优、1平、9劣，平均遗漏差为8.4。

SLS在C类的可行数为9/9，在U类为3/9。C类结果应结合其明确的簇间分离结构解释，不能据此推断空间构造在所有分布上都占优。

MIP有2次正常返回最优状态。全部方法中有1次被外部时限强制停止，其中0次保留了此前保存的可行方案；程序异常为0次。单次实测最大耗时为600.14秒。
